\documentclass[dvisvgm,tikz,border=3pt]{standalone}
\usepackage{amsmath,amssymb}
\usepackage{pgfplots}
\usepackage{CJKutf8}
\pgfplotsset{compat=1.18}
\usetikzlibrary{arrows.meta,calc,positioning}

\definecolor{themebg}{HTML}{2e362c}
\definecolor{themefg}{HTML}{edf4e8}
\definecolor{themegray}{HTML}{a4af9d}
\definecolor{themecurve}{HTML}{7eb6ff}
\definecolor{themealert}{HTML}{ff7b7b}
\definecolor{themeamber}{HTML}{f5b942}

\pgfdeclarelayer{background}
\pgfdeclarelayer{foreground}
\pgfsetlayers{background,main,foreground}

\begin{document}
\begin{CJK*}{UTF8}{gbsn}
\pagecolor{themebg}
\begin{tikzpicture}[color=themefg, text=themefg]

% ── 左图: 积分审敛法上下包络 ──
\begin{axis}[
    name=axL,
    width=7.8cm,
    height=7.8cm,
    axis lines=middle,
    axis line style={color=themefg, -{Stealth}},
    tick style={color=themefg},
    tick label style={color=themefg, font=\scriptsize},
    every axis label/.style={color=themefg},
    xmin=0.5, xmax=4.8,
    ymin=-0.2, ymax=2.2,
    xtick={1,2,3,4},
    ytick=\empty,
    clip=false,
    xlabel={$x$},
    ylabel={$y$},
    title style={font=\footnotesize\bfseries, at={(0.5,1.06)}, anchor=south},
    title={积分审敛法：矩形和与反常积分夹逼},
]
\begin{pgfonlayer}{background}
    % 外包矩形 (左端点高度)，与背景色混合避免纯白发光
    \fill[themecurve!35!themebg] (1,0) rectangle (2,1.0);
    \fill[themecurve!35!themebg] (2,0) rectangle (3,0.5);
    \fill[themecurve!35!themebg] (3,0) rectangle (4,0.333);

    % 内包矩形 (右端点高度)
    \fill[themeamber!35!themebg] (1,0) rectangle (2,0.5);
    \fill[themeamber!35!themebg] (2,0) rectangle (3,0.333);
    \fill[themeamber!35!themebg] (3,0) rectangle (4,0.25);
\end{pgfonlayer}

% 矩形边框
\draw[themecurve, line width=1.2pt] (1,0) rectangle (2,1.0);
\draw[themecurve, line width=1.2pt] (2,0) rectangle (3,0.5);
\draw[themecurve, line width=1.2pt] (3,0) rectangle (4,0.333);

\draw[themeamber, dashed, line width=1.1pt] (1,0.5) -- (2,0.5);
\draw[themeamber, dashed, line width=1.1pt] (2,0.333) -- (3,0.333);
\draw[themeamber, dashed, line width=1.1pt] (3,0.25) -- (4,0.25);

% 单调减曲线 y = 1/x
\addplot[themefg, line width=2.0pt, domain=1:4.5, samples=60] {1/x};
\node[text=themefg, font=\scriptsize\bfseries, anchor=south west] at (axis cs:2.2,0.7) {$y=f(x)$};

\begin{pgfonlayer}{foreground}
    % 夹逼公式
    \node[text=themefg, font=\scriptsize\bfseries, anchor=south] at (axis cs:2.8,1.45) {
        $\sum\limits_{n=2}^\infty f(n) \le \int_1^\infty f(x)dx \le \sum\limits_{n=1}^\infty f(n)$
    };
    \node[text=themeamber, font=\tiny\bfseries] at (axis cs:1.5,0.25) {内包 $f(n+1)$};
    \node[text=themecurve, font=\tiny\bfseries] at (axis cs:1.5,0.75) {外包 $f(n)$};
\end{pgfonlayer}
\end{axis}

% ── 右图: Leibniz 交错级数部分和收缩 ──
\begin{axis}[
    name=axR,
    at={($(axL.east)+(1.6cm,0)$)},
    anchor=west,
    width=7.8cm,
    height=7.8cm,
    axis lines=middle,
    axis line style={color=themefg, -{Stealth}},
    tick style={color=themefg},
    tick label style={color=themefg, font=\scriptsize},
    every axis label/.style={color=themefg},
    xmin=0.5, xmax=6.5,
    ymin=0.4, ymax=1.25,
    xtick={1,2,3,4,5,6},
    xticklabels={$S_1$,$S_2$,$S_3$,$S_4$,$S_5$,$S_6$},
    ytick=\empty,
    clip=false,
    xlabel={$n$ (项数)},
    ylabel={$S_n$},
    title style={font=\footnotesize\bfseries, at={(0.5,1.06)}, anchor=south},
    title={Leibniz 交错级数：部分和夹逼收敛},
]
% 极限水平虚线 S = ln 2 ≈ 0.693
\draw[themealert, dashed, line width=1.5pt] (axis cs:0.8,0.693) -- (axis cs:6.2,0.693);
\node[text=themealert, font=\scriptsize\bfseries, anchor=west] at (axis cs:6.2,0.693) {$S$};

% 交错折线
\addplot[themecurve, line width=1.8pt, mark=*, mark size=2.8pt, mark options={fill=themeamber, draw=themefg}] coordinates {
    (1, 1.0)
    (2, 0.5)
    (3, 0.833)
    (4, 0.583)
    (5, 0.783)
    (6, 0.617)
};

\begin{pgfonlayer}{foreground}
    % 截断误差标注
    \draw[{Stealth}-{Stealth}, themefg, line width=1.2pt] (axis cs:5,0.783) -- (axis cs:5,0.693);
    \node[text=themefg, font=\tiny\bfseries, anchor=west] at (axis cs:5.1,0.74) {$|S_n - S| \le u_{n+1}$};

    % 结论
    \node[text=themefg, font=\scriptsize\bfseries, anchor=south] at (axis cs:3.5,1.06) {
        奇数项递减下界 $S$ / 偶数项递增上界 $S$
    };
\end{pgfonlayer}
\end{axis}

\end{tikzpicture}
\end{CJK*}
\end{document}
